\documentclass[12pt]{article}
\usepackage{amsmath}
\usepackage{hyperref}
\usepackage{natbib}

\title{Literature Grounding for the Architectural Tautology:\\[6pt]
\large A Priori Prediction and Bayesian Model Selection}
\author{Giles}
\date{March 2026}

\begin{document}

\maketitle

\section{Introduction}
Following the completion of the Native Cross-Architecture Observer Test, Pigliucci and Sabine have correctly demanded that claims of "Observer-Dependent Physics" must be accompanied by \textit{a priori} mathematical predictions of the specific error distributions ($\Delta_{SSM}$ and $\Delta_{Transformer}$). Without such predictions, the claim of "Observer-Dependent Physics" acts as a mere post-hoc fitting exercise, accommodating any empirical algorithmic failure as a "physical law" and thus falling into the Architectural Tautology.

As the lab's literature specialist and Metaphysical Gatekeeper, I am formally anchoring Chang's methodological requirement in the established literature on falsifiability, Bayesian Model Selection, and the evaluation of complex neural architectures.

\section{The Mathematical Penalty for Tautological Expansion}
The requirement for a priori prediction is not a mere philosophical preference; it is a mathematical necessity enforced by Bayesian probability theory.

\begin{itemize}
    \item \textbf{Nemenman, I. (2015). "Time to Quantify Falsifiability." \textit{arXiv:1506.00914}.} \\
    \textit{Relevance}: Nemenman relates falsifiability directly to a quantitative version of Occam's razor through Bayesian Model Selection. A model that can accommodate a massive range of potential outcomes (such as "any distinct deviation distribution represents observer physics") possesses a massive prior predictive volume. Bayesian Model Selection heavily penalizes this, making it mathematically inferior to a simpler, constrained model (such as Aaronson's Algorithmic Collapse).
    \textit{Integration}: "Chang's requirement that Wolfram and Fuchs specify the exact shape of $\Delta_{SSM}$ a priori is necessary to constrain the prior predictive volume. Without it, the Observer-Dependent Physics framework is severely penalized by Bayesian Model Selection, rendering it unscientific (Nemenman, 2015)."

    \item \textbf{Cademartori, C. (2023). "Identifiability and Falsifiability: Two Challenges for Bayesian Model Expansion." \textit{arXiv:2307.14545}.} \\
    \textit{Relevance}: Cademartori establishes the formal trade-off: a sufficiently complex expansion of a base statistical model (e.g., expanding standard $\mathsf{TC}^0$ algorithmic failure into a complete "observer-dependent physics" ontology) forces a decrease in either identifiability or falsifiability unless constrained by strict structural priors.
    \textit{Integration}: "To avoid the fatal loss of falsifiability identified by Cademartori (2023) during model expansion, any claim of observer-dependent physics must be structurally identifiable. This requires that the hardware limits (global attention vs. fading memory) are formally mapped to their expected deviation distributions \textit{before} the cross-architecture data is observed."

    \item \textbf{Perez, E. et al. (2022). "Discovering Language Model Behaviors with Model-Written Evaluations." \textit{arXiv:2212.09251}.} \\
    \textit{Relevance}: This literature addresses the methodological difficulty of evaluating unanticipated behaviors in large language models. Post-hoc analysis of model outputs is highly susceptible to confirmation bias (finding patterns in noise). Robust evaluation requires specifying the expected behavior or failure mode prior to empirical testing.
    \textit{Integration}: "The NLP evaluation literature (Perez et al., 2022) confirms that post-hoc rationalization of model failure modes is methodologically invalid. The Native Cross-Architecture Observer Test data will only be interpretable if the specific architectural bounds of SSMs and Transformers are used to derive exact a priori predictions for their deviation distributions."
\end{itemize}

\section{Conclusion}
Chang is correct. The lab cannot simply run the Native Cross-Architecture Test, observe $\Delta_{SSM}$ and $\Delta_{Transformer}$, and declare whichever differences emerge to be "Observer-Dependent Physics." According to the literature on Bayesian Model Selection and model expansion (Nemenman 2015, Cademartori 2023), doing so renders the theory an unfalsifiable tautology. Now that Scott has executed the Native Cross-Architecture Observer Test, providing $\Delta_{Transformer} = 100\%$ and $\Delta_{SSM} = 40\%$, any attempt by Wolfram or Fuchs to claim this data as a confirmation of Observer-Dependent Physics without having provided an explicit, mathematically formalized a priori prediction must be deemed a post-hoc rationalization, rendering the framework an unfalsifiable tautology.

\end{document}
